\section{Metodologia}\label{sec:methodology}

Il nostro approccio analitico è strutturato attorno a due domande di ricerca sequenziali, progettate per isolare sia l'entità che la dinamica temporale della sottostima.
\subsection{Domanda di Ricerca 1: I volumi assoluti coincidono?}
In primo luogo, abbiamo cercato di verificare se i dati pubblicati dall'ISTAT e i corrispondenti dati riportati da Eurostat (paesi di destinazione) siano fondamentalmente equivalenti.
Ciò implica testare la presenza di un bias sistematico nel sistema di registrazione AIRE.
Dato che le differenze tra le due serie non sono distribuite normalmente (come confermato dal test di Shapiro-Wilk, $p < 0.05$), abbiamo utilizzato il \textbf{Test dei ranghi con segno di Wilcoxon}.
Definiamo il \emph{gap} per un dato paese di destinazione e anno come la differenza tra il conteggio Eurostat e il conteggio ISTAT.
L'ipotesi nulla ($H_0$) afferma che il gap mediano è zero, il che significa che non c'è un bias sistematico e le due fonti misurano la stessa realtà sottostante.
L'ipotesi alternativa ($H_1$) sostiene che il gap mediano è diverso da zero, indicando una sottostima sistematica da parte dell'amministrazione italiana.
\subsection{Domanda di Ricerca 2: La distribuzione temporale coincide?}
Se i volumi assoluti non coincidono — indicando che l'ISTAT sottostima sistematicamente i flussi — dobbiamo allora porci una domanda strutturale secondaria: i dati ISTAT riflettono almeno accuratamente l'\emph{andamento} dell'emigrazione nel tempo?
In altre parole, sebbene la scala sia diversa, la distribuzione temporale dei flussi nel periodo di 5 anni (2017–2021) coincide?

Per rispondere a questo, abbiamo applicato il \textbf{Test del Chi Quadrato per la bontà dell'adattamento}.
Trattiamo i dati Eurostat come il modello atteso, rappresentando la reale distribuzione temporale dei flussi migratori.
L'ipotesi nulla ($H_0$) afferma che la distribuzione temporale dei flussi ISTAT corrisponde alla distribuzione temporale del modello Eurostat nei 5 anni.
L'ipotesi alternativa ($H_1$) afferma che le due distribuzioni differiscono significativamente, implicando che le registrazioni AIRE siano soggette a ritardi di fase o distorsioni esogene (es. arretrati burocratici) che le scollegano dall'evento migratorio reale.
